\documentclass[11pt]{article}
\usepackage[a4paper,margin=25mm]{geometry}
\usepackage{amsmath,amssymb,amsthm,microtype,booktabs}
\usepackage[colorlinks=true,linkcolor=blue,urlcolor=blue,citecolor=blue]{hyperref}
\newtheorem{theorem}{Theorem}
\newtheorem{lemma}[theorem]{Lemma}
\newtheorem{corollary}[theorem]{Corollary}
\newcommand{\mis}{\operatorname{mis}}
\newcommand{\LC}{\operatorname{LC}}
\title{A square-root bound for global binomial log-concavity}
\author{}
\date{1 October 2026}
\begin{document}
\maketitle
\begin{abstract}
Let $\mathcal F\subseteq2^{[N]}$ contain the empty set, and let $M$ be
its number of inclusion-maximal members. We prove that every integer
$m\ge4(N+1)^4\sqrt M$ makes all coefficients of
$(1+x)^m\sum_{S\in\mathcal F}x^{|S|}$ log-concave.
No closure assumption on $\mathcal F$ is needed. The exponent $1/2$ of
$M$ cannot be reduced, even for independence complexes of trees, if
arbitrary fixed polynomial factors in $N$ are allowed.
Consequently the worst global repair thresholds are
$2^{N/4+o(N)}$ for forests, $3^{N/6+o(N)}$ for graphs, and
$2^{N/2+o(N)}$ for simplicial complexes or arbitrary set families.
The proof separates a coefficient-independent interior estimate from
uniform endpoint estimates. All constants in the upper bound are explicit.
\end{abstract}

\section{Statement and scope}
For a nonzero polynomial $P$ with nonnegative coefficients and positive
constant coefficient, put
\[
 \sigma(P)=\min\{m\in\mathbb Z_{\ge0}:
        (1+x)^mP(x)\text{ has a log-concave coefficient sequence}\}.
\]
Log-concavity here includes the absence of internal zeros, as well as
$c_k^2\ge c_{k-1}c_{k+1}$.
The assertions below prove finiteness in the cases under consideration.
For $\mathcal F\subseteq2^{[N]}$ containing $\varnothing$, write
$P_{\mathcal F}(x)=\sum_{S\in\mathcal F}x^{|S|}$, and let
$M(\mathcal F)$ count its inclusion-maximal members.

\begin{theorem}[Global square-root bound]\label{thm:main}
For $N\ge1$ and every such family $\mathcal F$,
\[
 \boxed{\quad\sigma(P_{\mathcal F})\le
       \left\lceil4(N+1)^4\sqrt{M(\mathcal F)}\right\rceil.\quad}
\]
In fact every integer at or above the displayed real bound makes the
entire coefficient sequence strictly log-concave at its interior indices.
\end{theorem}

For a graph $G$, use its family of independent sets. Then $M(\mathcal F)
=\mis(G)$ and multiplying by $(1+x)^m$ adds $m$ isolated vertices.
This theorem repairs \emph{all} log-concavity inequalities, not just the
last one. For a simplicial complex the same operation is joining with a
simplex on $m$ new vertices, or equivalently taking $m$ successive cones.

\paragraph{Position relative to existing work.}
Eventual log-concavity after multiplication by a binomial power is
classical: Handelman~\cite{H92,H11} treats a much wider existence question
and determines the exact threshold $k^2-3$ for $(1+x)^m(1+x^k)$.
The contribution derived here is an explicit bound governed by the
number of maximal members, with matching exponential rates for several
classes. The lower-bound tree construction is known from Galvin~\cite{G},
and the maximal-independent-set estimates are classical~\cite{PP}.
A literature search has not established publication priority for the
present quantitative statement. This is a self-contained research draft,
not an independently reviewed or formally verified result.

\section{The coefficient-independent interior}
Write $B_r=\binom mr$, with $B_r=0$ outside $0\le r\le m$.

\begin{lemma}[Pairwise binomial curvature]\label{lem:interior}
Let $n\ge1$, $L=n^2$, and $m\ge2L$. If
$L\le r,s\le m-L$ and $|r-s|\le n$, then
\[
 2B_rB_s>B_{r-1}B_{s+1}+B_{s-1}B_{r+1}.
\]
Consequently, for any $P(x)=\sum_{i=0}^{d}a_ix^i$ with $a_i\ge0$,
$a_0>0$, and $d\le n$, the coefficients
$c_k=\sum_i a_iB_{k-i}$ satisfy $c_k^2>c_{k-1}c_{k+1}$ whenever
\[
 n^2+n<k<m+d-(n^2+n).
\]
\end{lemma}
\begin{proof}
Set $u_r=r/(m-r+1)$ and $v_r=(m-r)/(r+1)$. Direct subtraction gives
\begin{align*}
 2-u_rv_s-u_sv_r
 &=\frac{(m+1)\{(r+1)(m-r+1)+(s+1)(m-s+1)
                  -(m+1)(r-s)^2\}}
 {(r+1)(s+1)(m-r+1)(m-s+1)}.
\end{align*}
The braces are strictly positive: their first two terms sum to at least
$2(L+1)(m-L+1)\ge(L+1)(m+2)>L(m+1)$, while the subtracted term is at most
$L(m+1)$. Multiply by $B_rB_s>0$.

For the consequence, every $r=k-i$ belongs to $[L,m-L]$ and any two
such values differ by at most $d\le n$. Symmetrize the identity
\[
 c_k^2-c_{k-1}c_{k+1}
 =\frac12\sum_{i,j}a_ia_j
 \{2B_{k-i}B_{k-j}-B_{k-1-i}B_{k+1-j}-B_{k-1-j}B_{k+1-i}\}.
\]
Every summand is nonnegative, and the $(0,0)$ summand is positive.
\end{proof}

\section{Uniform control at both endpoints}
\begin{lemma}[Two-endpoint smoothing]\label{lem:polynomial}
Let $n\ge1$, $d\le n$, $M\ge1$, and
$P(x)=\sum_{i=0}^{d}a_ix^i$ with $a_i\ge0$ and $a_0a_d>0$.
Suppose, for $2\le i\le d$, that
\[
 \frac{a_i}{a_0}\le M(n+1)^i,
 \qquad
 \frac{a_{d-i}}{a_d}\le M(n+1)^i.
\]
Then every integer $m\ge4(n+1)^4\sqrt M$ makes $(1+x)^mP$ strictly
log-concave at every interior index.
\end{lemma}
\begin{proof}
If $d=0$, this is the binomial sequence. Assume $d\ge1$, and put
$K=n(n+1)$. On the left, split
\[
 c_j=b_j+e_j,\qquad
 b_j=a_0B_j+a_1B_{j-1},\qquad
 e_j=\sum_{i=2}^{d}a_iB_{j-i}.
\]
For $0\le j\le K+1$ and $i\le j$,
\[
 \frac{B_{j-i}}{B_j}
 =\prod_{\ell=0}^{i-1}\frac{j-\ell}{m-j+\ell+1}
 \le\left(\frac{K+1}{m-K}\right)^i.
\]
Terms with $i>j$ vanish. Set $z=(n+1)(K+1)/(m-K)$.
As $K\le(n+1)^4\sqrt M$ and $K+1\le(n+1)^2$,
\[
 z\le\frac{1}{3(n+1)\sqrt M}\le\frac16,
 \qquad
 0\le\frac{e_j}{b_j}\le\frac{e_j}{a_0B_j}
 \le M\sum_{i=2}^{d}z^i\le2Mz^2
 \le\frac{2}{9(n+1)^2}\le\frac{1}{4K}.
\tag{1}
\]
All denominators are positive; the bound on $m$ also ensures $K+1<m$.

The baseline satisfies
\[
 b_k^2\ge\frac{k+1}{k}\,b_{k-1}b_{k+1}\qquad(1\le k\le K).
\tag{2}
\]
For completeness, write $u_j=j!b_j$. For $j\ge1$,
\[
 u_j=(m)_{j-1}L_j,\qquad
 L_j=(m-j+1)a_0+ja_1,
\]
where $(m)_r$ is a falling factorial. The $L_j$ are positive and affine,
so $L_j^2-L_{j-1}L_{j+1}=(a_1-a_0)^2\ge0$.
The falling-factorial factors are log-concave as well. This proves
$u_j^2\ge u_{j-1}u_{j+1}$ for $j\ge2$ in the required range.
At $j=1$, direct calculation gives $u_1^2-u_0u_2=ma_0^2+a_1^2\ge0$.
This is precisely (2).

With $\varepsilon=1/(4K)$, (1)--(2) imply, for $1\le k\le K$,
\[
 c_k^2\ge b_k^2\ge(1+1/K)b_{k-1}b_{k+1}
 >(1+\varepsilon)^2b_{k-1}b_{k+1}
 \ge c_{k-1}c_{k+1},
\]
since $1+1/K>(1+1/(4K))^2$.
Apply the identical argument to $x^dP(1/x)$ for the last $K$ indices
of the product, whose degree is $m+d$. All remaining indices satisfy
Lemma~\ref{lem:interior}. Positivity of all coefficients follows from
$a_0a_d>0$ and $m\ge d$: the contributions of $a_0$ and $a_d$ cover
$[0,m]$ and $[d,m+d]$, respectively.
\end{proof}

\section{From maximal members to all coefficients}
\begin{proof}[Proof of Theorem~\ref{thm:main}]
Put $d=\max_{S\in\mathcal F}|S|$ and $P_{\mathcal F}=\sum a_ix^i$.
We have $a_0=1$, $a_d\ge1$, and $a_i\le\binom Ni\le(N+1)^i$.
Every member of size $d-i$ extends to an inclusion-maximal member.
The extension adds at most $i$ elements. Each maximal member contains at
most $\sum_{j=0}^{i}\binom Nj$ subsets obtained by deleting at most $i$
elements. Therefore
\[
 a_{d-i}\le M\sum_{j=0}^{i}\binom Nj\le M(N+1)^i.
\tag{3}
\]
The last inequality follows by encoding a subset of size at most $i$ as
its increasing list, padded with zeros to length $i$, over an alphabet
of $N+1$ symbols. This encoding is injective. Division by $a_d\ge1$
only improves the estimate. Lemma~\ref{lem:polynomial}, with $n=N$,
now applies. The case $d=0$ is already included.
\end{proof}

\begin{corollary}[A polynomial bound for few maximal members]
If $M(\mathcal F)\le N^c$ for a fixed $c$, then
$\sigma(P_{\mathcal F})\le\lceil4(N+1)^4N^{c/2}\rceil$.
This applies without assuming that the family is a simplicial complex.
\end{corollary}

\section{Sharp exponential rates}
Let $S_{\rm for}(N)$ and $S_{\rm gr}(N)$ be the maximum of
$\sigma(I(G;x))$ over forests and all graphs, respectively, of order at
most $N$. Let $S_{\rm cx}(N)$ and $S_{\rm fam}(N)$ be the analogous
maxima over simplicial complexes and set families containing the empty
set, on at most $N$ labelled elements.

\begin{theorem}\label{thm:rates}
As $N\to\infty$,
\[
 \boxed{\begin{aligned}
 S_{\rm for}(N)&=2^{N/4+o(N)},\\
 S_{\rm gr}(N)&=3^{N/6+o(N)},\\
 S_{\rm cx}(N)&=2^{N/2+o(N)},\\
 S_{\rm fam}(N)&=2^{N/2+o(N)}.
 \end{aligned}}
\]
These are thresholds for global log-concavity.
\end{theorem}
\begin{proof}
We give both sides.

\smallskip\noindent\textit{Upper bounds.}
The classical bounds $\mis(G)\le2^{N/2}$ for forests and
$\mis(G)\le3^{N/3}$ for graphs give
\[
 S_{\rm for}(N)\le\lceil4(N+1)^4 2^{N/4}\rceil,\qquad
 S_{\rm gr}(N)\le\lceil4(N+1)^4 3^{N/6}\rceil.
\tag{4}
\]
Here are short proofs of the counting inputs. In a forest, an isolated
vertex is forced into every maximal independent set. Otherwise choose
a leaf $v$ with neighbor $u$; partitioning maximal sets by which of
$u,v$ they contain gives
\[
 \mis(G)=\mis(G-\{u,v\})+\mis(G-N[u])\le2\cdot2^{(N-2)/2}.
\]
For a general graph without isolates, choose a minimum-degree vertex
$v$, of degree $\delta\ge1$. Every maximal independent set meets $N[v]$;
conditioning on a chosen member and overcounting gives
\[
 \mis(G)\le\sum_{u\in N[v]}\mis(G-N[u])
 \le(\delta+1)3^{(N-\delta-1)/3}\le3^{N/3}.
\]
The last inequality is $k\le3^{k/3}$ for integers $k\ge2$, proved from
$k=2,3,4$ and the step $k+3\le3k$. These arguments are inductions on $N$.

The maximal members of any set family form an antichain. Counting
permutations whose initial segment is a member of an antichain gives
$\sum_{S\text{ maximal}}\binom N{|S|}^{-1}\le1$. Thus, writing
$B_N=\binom N{\lfloor N/2\rfloor}$,
\[
 S_{\rm fam}(N)\le\lceil4(N+1)^4\sqrt{B_N}\rceil,
\tag{5}
\]
and the same bound holds for complexes.

\smallskip\noindent\textit{Graph lower bounds.}
For $q\ge2$ and $t\ge1$, let $G_{q,t}$ have a root with three neighbors
$w_1,w_2,w_3$. Attach $t$ disjoint $q$-cliques to each $w_i$, joining
$w_i$ to exactly $q-1$ vertices of every attached clique. Then
\[
 |V(G_{q,t})|=3qt+4,\quad \alpha=3t+3,\qquad
 I(G_{q,t};x)=\big((1+qx)^t+x(1+x)^t\big)^3+x(1+qx)^{3t}.
\]
For $q=2$ this is Galvin's known tree $T_{3,t}$~\cite{G}.
Its top three coefficients, in descending order, are $1,A,B$, where
\[
 A=3(q^t+t),\qquad
 B=q^{3t}+3\left(tq^{t-1}+\binom t2\right)+3(q^t+t)^2.
\]
After $m$ isolates, the last log-concavity difference is
\[
 A^2-B+Am+\frac{m(m+1)}2.
\tag{6}
\]
For fixed $q$, $B-A^2\sim q^{3t}$ and $A=o(q^{3t/2})$.
The exact positive root of (6) consequently gives the terminal threshold
$\tau(G_{q,t})\sim\sqrt2\,q^{3t/2}$. Since $\sigma\ge\tau$, the
choices $q=2,3$ match (4) in exponential rate.
To cover every $N$, choose $t=\lfloor(N-4)/(3q)\rfloor$.

An explicit finite lower bound is also available: for $q\ge2,t\ge8$,
$A\le4q^t$ and $B-A^2\ge15q^{3t}/16$. At
$m=\lfloor q^{3t/2}\rfloor$, we have
$Am\le q^{3t}/4$ and $m(m+1)/2\le9q^{3t}/16$.
Hence
\[
 S_{\rm for}(N)>2^{N/4-5/2}\ (N\ge52),\qquad
 S_{\rm gr}(N)>3^{N/6-13/6}\ (N\ge76).
\tag{7}
\]

\smallskip\noindent\textit{Complex lower bounds.}
For $N\ge4$, put $r=\lfloor N/2\rfloor$, $d=r+2$, and fix a $d$-set
$A\subseteq[N]$. The simplicial complex
\[
 \mathcal K_N=\{S\subseteq[N]:|S|\le r\}\ \cup\ 2^A
\]
has polynomial
\[
 P_{\mathcal K_N}(x)=\sum_{j=0}^{r}\binom Njx^j+dx^{d-1}+x^d.
\]
Its top coefficients are $1,d,B_N$. By (6), its terminal threshold is
$\sim\sqrt{2B_N}=2^{N/2+o(N)}$, because $d=O(N)$ and
$2^N/(N+1)\le B_N\le2^N$. This proves both remaining lower bounds.
For a finite version, $N\ge16$ implies $B_N\ge4d^2$:
$2^N\ge4(N+1)(N+2)^2$ follows by induction from $N=16$.
At $m=\lfloor\sqrt{B_N}/2\rfloor$, the positive terms in (6) sum to
at most $B_N/2$, while $B_N-d^2\ge3B_N/4$. Thus
\[
 S_{\rm cx}(N)>\tfrac12\sqrt{B_N}
 \ge\frac{2^{N/2}}{2\sqrt{N+1}}\qquad(N\ge16).
\]
\end{proof}

\paragraph{Sharpness in the maximal-member parameter.}
For the same graphs, put $x=q^t$. If the root is selected, there are
$x^3$ maximal independent sets. If it is not selected, at least one
$w_i$ must be selected. In each branch, selecting $w_i$ gives one maximal
configuration; omitting it gives $x-1$, since at least one selected clique
vertex must dominate $w_i$. Thus
\[
 \mis(G_{q,t})=x^3+\{x^3-(x-1)^3\}
             =x^3+3x^2-3x+1\sim q^{3t}.
\]
Therefore $\sigma(I(G_{q,t};x))\ge(\sqrt2+o(1))\sqrt{\mis(G_{q,t})}$.
No bound $N^C M^{1/2-\varepsilon}$, with fixed $C$ and
$\varepsilon>0$, can hold for all trees. The polynomial factor $4(N+1)^4$
in our upper bound is not claimed optimal.

\section{Verification and limits of the claim}
The accompanying exact-arithmetic verifier passed 191,941 checks, including
all 138 families containing the empty set on one, two, or three elements.
It checks the pairwise identity,
its prescribed interior range, the baseline curvature, the endpoint
constants, finite set-family counting bounds, and concrete global
thresholds. These are checks supporting the written universal proof;
they are not a proof-assistant certificate. In particular, a finite
experiment is not used to infer either asymptotic theorem.

The result does not prove that the global and terminal thresholds agree
for each graph. It proves that their \emph{worst-case exponential rates}
agree in the graph classes above. It does not resolve the number of
log-concavity failures of an unmodified forest. The present argument and
its audit were developed with ChatGPT and have not undergone independent
mathematical review.

\begin{thebibliography}{9}
\bibitem{H92} D. Handelman, \emph{Spectral radii of primitive integral
companion matrices and log concave polynomials}, Contemporary Mathematics
135 (1992), 231--237.
\bibitem{H11} D. Handelman, \emph{Log concavity of $(1+x)^m(1+x^k)$},
arXiv:1102.2961, version 2 (2018).
\url{https://arxiv.org/abs/1102.2961}.
\bibitem{G} D. Galvin, \emph{Trees with non log-concave independent set
sequences}, arXiv:2502.10654 (2025).
\url{https://arxiv.org/abs/2502.10654}.
\bibitem{PP} C. Palmer and B. Patk\'os, \emph{On the number of maximal
independent sets: From Moon--Moser to Hujter--Tuza}, Journal of Graph
Theory 102 (2023), 440--445; arXiv:2205.04082.
\url{https://arxiv.org/abs/2205.04082}.
\end{thebibliography}
\end{document}
